\documentclass[12pt]{article}

\usepackage[margin=1in]{geometry}
\usepackage{amsmath,amsthm,amssymb}
\usepackage[all]{xy}
  \SelectTips{cm}{10}
  \everyxy={<2.5em,0em>:}
\DeclareMathOperator{\Ima}{Im}

\newcommand{\Z}{\mathbb{Z}}
\newcommand{\lean}[1]{\texttt{#1}}

\newenvironment{ex}[2][Problem]{\begin{trivlist}
\item[\hskip \labelsep {\bfseries #1}\hskip \labelsep {\bfseries #2.}]}{\end{trivlist}}

\newenvironment{sol}[1][Solution]{\begin{trivlist}
\item[\hskip \labelsep {\bfseries #1:}]}{\end{trivlist}}


\begin{document}

% --------------------------------------------------------------
%                         Start here
% --------------------------------------------------------------

\begin{ex}{1} [Problem 9 from HW01.lean]

    Show that there exists a proposition $Q$ such that for any proposition $P$, we have that
    $P\vee Q$ is true if and only if $P$ is true.

\end{ex}


\begin{proof}

    I claim that the proposition \lean{False} satisfies this property. Let $P$ be an arbitrary
    proposition. Consider the proposition $P\vee\lean{False}$. Now we will explore the truth table
    verbally: if $P$ is true, then $P\vee\lean{False}$ is true; if $P$ is false, then
    $P\vee\lean{False}$ is false. Therefore $P\vee\lean{False}$ is true only when $P$ is true,
    so $P\vee\lean{False}$ implies $P$. Conversely, $P$ implies $P\vee\lean{False}$ by the definition
    of the logical connective "or." At last, we have proven implication both ways, therefore
    $P\vee\lean{False}$ is true if and only if $P$ is true. Since \lean{False} satisfies the desired
    property for the arbitrarily chosen proposition $P$, then \lean{False} satisfies the general
    property described of $Q$, so such a proposition $Q$ does exist.

\end{proof}

\begin{ex}{2} [Infinitely many primes]

    Recall that a natural number $p$ is called "prime" if and only if the only natural numbers
    dividing $p$ are $1$ and $p$ itself. 1 is not considered prime, since it has only 1 distinct
    factor.
    There are infinitely many natural numbers which are prime.

\end{ex}

\begin{proof}{} [Due to Euclid, with some tweaks]

    We will prove the claim by way of contradiction. Suppose there are finitely many prime
    numbers. Let $S$ denote that finite set of prime numbers. Now we construct a natural number
    $N$, defined as the product of all the elements in $S$, i.e. $N:=\Pi_{p\in S}p$.
    By construction, every prime number divides $N$. Since we know 1 has only itself as
    a divisor, then every prime number does not divide 1. Let $p$ be a prime number.
    Then $p\mid N$ and $p\nmid 1$, thus $p\nmid (N+1)$. Therefore $N+1$ has no prime divisors.
    By theorem, the only natural number with no prime divisors is 1, so $N+1=1$. 1 is the smallest
    natural number, so $N\ge1$, thus $N+1>1$. Then through substitution we have that $1>1$,
    which is plainly a contradiction. Since we have arrived at a contradiction, our most recent
    assumption---that there are finitely many primes---is false. Hence there are infinitely
    many natural numbers which are prime.


\end{proof}

\end{document}